\documentclass[11pt]{article}
\usepackage[margin=0.68in]{geometry}
\usepackage{booktabs,enumitem,fancyhdr,graphicx,overpic,tabularx,titlesec}
\usepackage[table]{xcolor}
\usepackage[hidelinks]{hyperref}
\usepackage{circuitikz}
\graphicspath{{docs/lessons/assets/}{../assets/}}
\hypersetup{pdftitle={ADK Lesson 039: Percussion Sequencer},
 pdfauthor={Kevin Shanahan},
 pdfsubject={Contact grouping, acoustic association, quantization, and playback intents},
 pdflang={en-US}}
\pagestyle{fancy}\fancyhf{}
\lhead{\textsc{ADK Lesson 039}}\rhead{Percussion Sequencer}
\cfoot{\thepage}\setlength{\headheight}{14pt}
\setlength{\parindent}{0pt}\setlength{\parskip}{5pt}
\setlist{nosep,leftmargin=1.45em}
\titleformat{\section}{\Large\bfseries}{\thesection}{0.5em}{}
\newcommand{\lessonbox}[1]{\begin{center}\fbox{\parbox{0.92\linewidth}{#1}}\end{center}}

\begin{document}
\begin{center}
 {\Huge\bfseries Percussion Sequencer}\\[4pt]
 {\Large Group attacks, associate one envelope, replay bounded cues}\\[8pt]
 \textit{Exact reference design; incoming conformance and E1 acceptance open}
\end{center}
\begin{tabularx}{\linewidth}{@{}>{\bfseries}l X >{\bfseries}l X@{}}
Lesson & 039 --- project & Time & 180--240 minutes\\
Prerequisites & Lessons 037--038 & Board & Mega 2560 reference fixture\\
Evidence & six LEDs, passive piezo, named test points &
Status & integration verification active; bench open\\
\end{tabularx}

\lessonbox{\textbf{Scope gate.} Four C\&K/Littelfuse
\texttt{RB-220-07A R} contacts and one SparkFun \texttt{SEN-12642 V10}
define the reference circuit. A look-alike kit part is not interchangeable.
Intensity is relative ADC evidence, not sound pressure, loudness, force,
damage, speech, or a safety measurement.}

\section{From four lanes to one pattern}
\begin{center}
% ADK visual: pencil
\begin{overpic}[width=0.96\linewidth]{039-pattern-pencil.png}
 \put(4,5){\small four qualified lanes}
 \put(49,5){\small quantized grid}
 \put(84,5){\small current frame}
 \put(52,36){\small supplied-time cursor}
\end{overpic}

\small\textit{Alternative text: four hand-drawn lanes feed a bounded
sixteenth-note grid. A supplied-time cursor selects only the current frame.}
\end{center}

The first qualified attack freezes recording tempo and epoch. Up to one attack
per lane enters an inclusive simultaneous window. One eligible completed
acoustic interval supplies relative intensity to every group member; timeout
supplies zero with explicit \texttt{AssociationTimeout} provenance. Only one
closed group waits, so later attacks are visibly suppressed rather than
secretly queued.

The pattern contains 4--16 sixteenth-note steps at 30--240 BPM:
\[
 \mathrm{stepPeriodMs}=\left\lfloor\frac{60000}{4\,\mathrm{tempoBpm}}\right\rfloor .
\]
Nearest-step quantization rounds an exact half-step forward. Fixed storage
holds 32 hits. A group that cannot fit is rejected atomically.
\texttt{hit(index)} returns a copied hit with its association provenance.

\section{Electrically authoritative input schematic}
The input schematic below is authoritative for this named reference fixture only.
Each contact has its own 10 k\(\Omega\) external pull-up. The adapter therefore
uses \texttt{Pull::None} and interprets closure as active low.

\begin{center}
% ADK visual: schematic
\begin{circuitikz}[american]
 \draw (0,4) node[vcc]{5 V} to[R,l=10 k$\Omega$] (0,2.8)
       node[circ]{} to[push button] (0,1.4) node[ground]{};
 \node[below left] at (0,2.8){TP-C0};
 \draw (2,4) node[vcc]{5 V} to[R,l=10 k$\Omega$] (2,2.8)
       node[circ]{} to[push button] (2,1.4) node[ground]{};
 \node[below left] at (2,2.8){TP-C1};
 \draw (4,4) node[vcc]{5 V} to[R,l=10 k$\Omega$] (4,2.8)
       node[circ]{} to[push button] (4,1.4) node[ground]{};
 \node[below left] at (4,2.8){TP-C2};
 \draw (6,4) node[vcc]{5 V} to[R,l=10 k$\Omega$] (6,2.8)
       node[circ]{} to[push button] (6,1.4) node[ground]{};
 \node[below left] at (6,2.8){TP-C3};
 \draw (9,4) rectangle (12,1.3);
 \node at (10.5,3.5){SEN-12642 V10};
 \draw (9,2.8) -- (8.1,2.8) node[left]{A8};
 \node[right] at (9,2.8){ENVELOPE};
 \draw (9,2.2) -- (8.1,2.2) node[left]{D44};
 \node[right] at (9,2.2){GATE};
 \node[right] at (9,1.65){AUDIO NC};
 \draw (11.3,4) -- (11.3,4.5) node[vcc]{5 V};
 \draw (11.3,1.3) -- (11.3,.8) node[ground]{};
\end{circuitikz}

\small Four nonpolar rolling-ball contacts close TP-C0--TP-C3 to ground.
The module and Mega share 5 V and ground; AUDIO is disconnected.
\end{center}

\begin{tabularx}{\linewidth}{@{}l l X@{}}
\toprule
Path & Mega pin & Observation\\
\midrule
TP-C0--TP-C3 & D40--D43 & high open; low when its reference contact closes\\
TP-E / TP-T & A8 / D44 & envelope ADC and active-high gate\\
TP-BPM & A9 & 10 k\(\Omega\) potentiometer wiper\\
surface evidence & D45--D48 & accepted/current surface, each through 1 k\(\Omega\)\\
step display & D49--D51 & 74HC595 data, clock, and latch\\
state/capture & D52/D53 & ready/heartbeat and hardware-sample interval\\
cue & D6 / timer 2 & identified passive piezo, bounded tone only\\
play/clear & D38/D39 & buttons to ground using internal pull-ups\\
\bottomrule
\end{tabularx}

\section{Observe, decide, actuate}
The Mega adapter first updates four inputs and \texttt{ContactDynamics}
policies, then the envelope and gate inputs and \texttt{AcousticEnvelope}.
It copies only qualified attack bits, source statuses, and a completed
acoustic record into \texttt{PercussionSequencerInput}. The engine owns no
pin, timer, clock, endpoint, or dynamic storage.

Playback position derives from elapsed supplied time. Sparse calls skip to the
frame current now and never replay crossed cues late. A one-digit display
shows the active step. Accepted-hit LEDs latch for 180 ms and show current
surfaces during playback. In reference-hardware mode only, D53 toggles on each
coalesced 20 ms hardware sample; replay leaves it inactive.
The strongest current hit selects 262, 330, 392, or 523 Hz; equal intensity
selects the lowest lane. A frame without a hit has zero tone intent. That does
not prove physical silence.

\section{Predict, observe, interpret}
\begin{enumerate}
 \item Keep \texttt{ADK\_LESSON039\_USE\_REFERENCE\_HARDWARE=0}.
       Predict three replay hits followed by playback from step zero.
 \item Observe D45--D48 for accepted/current surfaces. Jumping time must not
       produce a burst of missed LED or tone cues.
 \item Observe D52 steady as ready outside playback and changing once per beat
       during playback, including empty pattern regions. D53 stays inactive
       in replay; on a qualified reference build it toggles on each coalesced
       20 ms hardware sample.
 \item Inject one failing source status in host replay. Predict retained host
       fault attribution while D52, D53, the display, surface LEDs, and piezo
       all become inactive.
 \item At the simultaneous boundary predict inclusion; one tick later predict
       suppression while association waits.
 \item Omit completion. Predict zero intensity with timeout provenance, not a
       claim that the room was quiet.
\end{enumerate}

\section{Qualified-build and shutdown checklist}
Before selecting hardware mode, match all five incoming specimens to the
named references, inspect unpowered wiring, and approve a bounded powered
plan. Measure the actual 5 V rail and every pull-up; prove each worst-case
closed current is no greater than the RB reference's 1 mA maximum. Measure
ENVELOPE and GATE ranges against Mega limits. Complete per-pin, per-port,
74HC595-package, 5 V rail, and total-current arithmetic; the known direct LED
and display bound is 70 mA, excluding module, pulls, passive piezo, and Mega
baseline. Power the module and Mega together; never drive a module output into
an unpowered Mega.

Resource acquisition and electrical safe state are separate evidence.
Initialization status establishes claims. After shutdown, verify LEDs and
piezo inactive and separately inspect TP-C0--TP-C3, TP-E, and TP-T. Remove USB
for heat, odor, reset, unstable rail, or unexpected output. Rewire only while
unpowered.

\lessonbox{\textbf{Physical acceptance remains open.} Compilation, host
replay, documentation, and an authoritative reference schematic do not prove
the delivered specimens, currents, output ranges, closure orientation,
sounder behavior, or safe shutdown on a bench. Record those observations only
after the qualified fixture is measured.}

\section{Software evidence and sources}
The project-engine tests exercise grouping and association boundaries,
quantization, capacity, playback, source attribution, lifecycle, and the
versioned fieldwise trace. Focused example tests exist for replay scheduling,
frame-cue gating, and the hardware-free adapter path; their shared promotion
gate remains integration work.

\begin{itemize}
 \item C\&K/Littelfuse RB-series formal datasheet, exact
       \texttt{RB-220-07A R} drawing and ratings.
 \item SparkFun SEN-12642 hardware V1.0 source at repository commit
       \texttt{a456d9e}, product page, and hookup guide.
 \item Texas Instruments LMV324 primary datasheet.
 \item ADK sensor evidence report, testing contract, and safety model.
\end{itemize}

This untagged print edition has a searchable HTML alternative.
It makes no PDF/UA or physical accessibility claim.
\end{document}
